% ============================================================
% Section 97: Er₄Nb₄ 电阻率-温度关系分析
% ============================================================
\section{固体理论：Er$_4$Nb$_4$ 电阻率与温度的关系}
\label{sec:er4nb4-resistivity}

\begin{modelbox}[题目：掺杂金属电阻率的温度依赖与平均自由程]
一固体（Er$_4$Nb$_4$，掺杂，但并未有意损伤）的电阻率与温度的关系如图 1.79 所示。

(1) 从图中看，该固体是金属还是绝缘体？
(2) 描述产生电阻的主要物理过程，并解释在以下三个区内电阻率与温度的关系：
\quad(a) $T\sim0\,$K；\quad(b) $T\sim25\,$K；\quad(c) $T\sim300\,$K。
(3) 估计 $T=0\,$K 时和 $T=300\,$K 时的电子平均自由程和平均自由时间。此材料是好金属吗？

（已知：$n=10^{23}\,\text{cm}^{-3}$，$e=5\times10^{-10}\,\text{esu}$，$v_F=10^8\,\text{cm/s}$，$m_e=10^{-27}\,\text{g}$）
\end{modelbox}

\begin{tipbox}[考点快速分析]
\begin{itemize}
    \item \textbf{金属 vs 绝缘体}：电阻率随 $T$ 增大（正温度系数）为金属；指数下降为高阻绝缘体
    \item \textbf{马西森定则}：$\rho(T)=\rho_{\text{imp}}+\rho_{\text{ph}}(T)$
    \item \textbf{低温 $T^5$ 律}：布洛赫-格林艾森，小角声子散射
    \item \textbf{高温线性区}：$\rho\propto T$，大角声子散射
    \item \textbf{平均自由程}：$l=v_F\tau$，$\tau=m/(ne^2\rho)$（CGS单位）
\end{itemize}
\end{tipbox}

\begin{derivationbox}[标准解题过程]
\textbf{(1) 金属还是绝缘体？}

从图 1.79 可见，该材料电阻率随温度升高而增大（正温度系数），且室温下 $\rho\sim10^{-4}\,\Omega\cdot\text{cm}$。绝缘体则具有负温度系数且室温电阻率极高（$>10^{10}\,\Omega\cdot\text{cm}$）。

\begin{equation}
\boxed{\text{该固体是金属。}}
\end{equation}

\textbf{(2) 各区电阻率与温度的关系}

金属电阻来源于电子受晶格振动（声子）和晶体缺陷（杂质、空位等）的散射。根据\textbf{马西森定则}：
\begin{equation}
\rho(T)=\rho_{\text{imp}}+\rho_{\text{ph}}(T),
\end{equation}
其中 $\rho_{\text{imp}}$ 为与温度无关的剩余电阻，$\rho_{\text{ph}}(T)$ 为声子散射贡献的理想电阻。

\textbf{(a) $T\sim0\,$K（极低温区）}

声子热激发几乎为零，$\rho_{\text{ph}}\approx0$。电阻由杂质和缺陷散射主导，即剩余电阻 $\rho_{\text{imp}}$ 为常数。图中 $T\to0$ 时 $\rho\approx27\times10^{-6}\,\Omega\cdot\text{cm}$ 趋于水平。

\textbf{(b) $T\sim25\,$K（低温区）}

只有长波声子（小 $q$）被热激发。小 $q$ 声子散射导致电子动量改变的角度 $\theta\sim q/k_F$ 很小，一次碰撞不足以消除电场方向的定向动量。需要约 $(k_F/q)^2\propto T^{-2}$ 次碰撞才能构成一次有效散射（使动量完全弛豫）。

低温下声子数密度 $n_{\text{ph}}\propto T^3$（德拜模型）。有效散射率
\begin{equation}
\frac{1}{\tau}\propto n_{\text{ph}}\times\left(\frac{q}{k_F}\right)^2\propto T^3\times T^2=T^5.
\end{equation}

因此理想电阻率满足\textbf{布洛赫-格林艾森 $T^5$ 定律}：
\begin{equation}
\boxed{\rho_{\text{ph}}\propto T^5.}
\end{equation}

\textbf{(c) $T\sim300\,$K（高温区）}

大量声子被激发，$n_{\text{ph}}\propto T$。声子波矢与费米波矢可比，每次碰撞散射角大，几乎一次碰撞即可有效弛豫动量。散射率 $1/\tau\propto n_{\text{ph}}\propto T$。因此
\begin{equation}
\boxed{\rho\propto T,}
\end{equation}
呈现线性温度依赖关系。

\textbf{(3) 平均自由程与平均自由时间的估算}

在 CGS 单位制中，电导率 $\sigma=ne^2\tau/m$，故
\begin{equation}
\tau=\frac{m}{ne^2\rho},\qquad l=v_F\tau=\frac{mv_F}{ne^2\rho}.
\end{equation}

从图 1.79 读取：$T=0\,$K 时 $\rho\approx27\times10^{-6}\,\Omega\cdot\text{cm}$；$T=300\,$K 时 $\rho\approx120\times10^{-6}\,\Omega\cdot\text{cm}$。

代入题目给定数值：
\begin{align}
\tau(0\,\text{K})&=\frac{10^{-27}}{10^{23}\times(5\times10^{-10})^2\times27\times10^{-6}}
\approx1.5\times10^{-13}\,\text{s},\\
l(0\,\text{K})&=10^8\times1.5\times10^{-13}=1.5\times10^{-5}\,\text{cm}=\boxed{1500\,\text{\AA}.}
\end{align}

\begin{align}
\tau(300\,\text{K})&=\frac{10^{-27}}{10^{23}\times(5\times10^{-10})^2\times120\times10^{-6}}
\approx4.5\times10^{-14}\,\text{s},\\
l(300\,\text{K})&=10^8\times4.5\times10^{-14}=4.5\times10^{-6}\,\text{cm}=\boxed{45\,\text{\AA}.}
\end{align}

\textbf{是否是好金属？}

良好金属（如 Cu、Ag）室温下平均自由程通常为数百埃（Cu 约 400\,\AA）。该材料室温 $l\sim45\,\text{\AA}$，仅为原子间距的十余倍，电子散射较强。因此
\begin{equation}
\boxed{\text{此材料不是好金属。}}
\end{equation}
\end{derivationbox}

\begin{formulabox}[答案汇总]
\begin{center}
\begin{tabular}{ll}
\toprule
\textbf{问题} & \textbf{结论} \\
\midrule
(1) 金属/绝缘体 & 金属 \\
(2a) $T\sim0\,$K & 剩余电阻主导，$\rho\approx$ 常数 \\
(2b) $T\sim25\,$K & 理想电阻 $\rho_{\text{ph}}\propto T^5$（布洛赫-格林艾森律） \\
(2c) $T\sim300\,$K & 理想电阻 $\rho_{\text{ph}}\propto T$（高温线性区） \\
(3) $T=0\,$K & $\tau\approx1.5\times10^{-13}\,$s，$l\approx1500\,\text{\AA}$ \\
(3) $T=300\,$K & $\tau\approx4.5\times10^{-14}\,$s，$l\approx45\,\text{\AA}$ \\
(3) 是否好金属 & 否（室温平均自由程偏小） \\
\bottomrule
\end{tabular}
\end{center}
\end{formulabox}

\begin{insightbox}[物理图像]
\begin{itemize}
    \item 金属电阻的温度依赖性是电子-声子散射机制的宏观体现。低温 $T^5$ 律源于长波声子的小角散射特性：只有积累足够多的小角散射才能有效弛豫电子的定向动量。
    \item 剩余电阻反映了样品的纯度和晶体完整性，是低温下电阻的极限。马西森定则将两种独立的散射机制线性相加。
    \item 通过电阻率和费米速度估算的平均自由程，是衡量金属导电品质的微观参数。$l\sim45\,\text{\AA}$ 表明其电子输运受较强散射限制，可能与掺杂引入的杂质散射有关。
\end{itemize}
\end{insightbox}
